\begin{proofbox}
	\textbf{\textcolor{CProof}{思路提示}}

	球心向截面作垂线，垂足恰为截面圆心；连接球心和截面圆周上任一点，得到球半径，与垂线、截面半径组成直角三角形。

	\medskip
	\textbf{\textcolor{CProof}{正式推导}}
	\begin{description}[style=nextline, leftmargin=2.8em, labelsep=0.5em, itemsep=0.55em]
		\item[\textbf{步骤一（作辅助线）：}] 设球心$O$，截面圆心$O'$，$OO'$垂直于截面，且$OO'=d$。在截面上取一点$P$，连接$OP$、$O'P$。
			\par\textbf{理由：} $O'$是球心在截面上的投影，所以$OO'$垂直于截面内任何过$O'$的直线，包括$O'P$。因此$\triangle OO'P$是直角三角形，$\angle OO'P=90^\circ$。
		\item[\textbf{步骤二（标识已知量）：}] $OP=R$（球半径），$OO'=d$（球心到平面距离），$O'P=r$（截面圆半径）。
			\par\textbf{理由：} 这些是已知定义。
		\item[\textbf{步骤三（应用勾股定理）：}] 在直角三角形$OO'P$中，斜边是$OP$，直角边是$OO'$和$O'P$，由勾股定理：
			\[
				OP^{2} = OO'^{2} + O'P^{2} \Longrightarrow R^{2} = d^{2} + r^{2}.
			\]
			\par\textbf{理由：} 直角三角形两直角边的平方和等于斜边的平方。
		\item[\textbf{步骤四（解出截面半径）：}] 将$r$视为未知数，从$R^{2}=d^{2}+r^{2}$中解出$r$：
			\[
				r^{2} = R^{2} - d^{2} \Longrightarrow r = \sqrt{R^{2} - d^{2}}.
			\]
			\par\textbf{理由：} 若$0\le d<R$，则$R^{2}-d^{2}>0$，开平方取正值（半径非负）。
	\end{description}

	\medskip
	\textbf{\textcolor{CProof}{结论回扣}}

	\textbf{因此，球截面圆半径公式$r=\sqrt{R^{2}-d^{2}}$成立，其中条件$0\le d<R$保证根号内为正，且$r\ge0$。该公式核心来自直角三角形勾股定理。}
\end{proofbox}
